\SourceLeaf{176}Es ist leicht zu beweisen, dass, wenn Euklids Geometrie nicht die wahre ist,
es gar keine ähnliche Figuren gibt: die Winkel in einem gleichseitigen Dreieck
sind dann auch nach der Grösse der Seite verschieden, wobei ich gar nichts
absurdes finde. Es ist dann der Winkel Function der Seite und die Seite Function
des Winkels, natürlicher Weise eine solche Function, in der zugleich eine
constante Linie vorkommt. Es scheint etwas paradox, dass eine constante Linie
  gleichsam a priori möglich sein könne; ich finde aber darin nichts
widersprechendes. Es wäre sogar wünschenswerth, dass die Geometrie Euklids nicht
  wahr wäre, weil wir dann ein allgemeines Mass a priori hätten, z. B.
könnte man als Raumeinheit die Seite desjenigen gleichseitigen Dreiecks annehmen,
  dessen Winkel $=59^{0} 59' 59''{,}99999$. .\ .\ .\ .\ .\ .\ .

\begin{center}\rule{3cm}{0.4pt}\end{center}
